% SKY2100 -- Exercise 8 (self-contained article; no external .sty needed)
\documentclass[11pt]{article}

% ---- Packages ------------------------------------------------------
\usepackage[a4paper,margin=2.2cm]{geometry}
\usepackage[T1]{fontenc}
\usepackage[utf8]{inputenc}
\usepackage{xcolor}
\usepackage{enumitem}
\usepackage{pgffor}
\usepackage{amssymb}
\usepackage{array}

% ---- Colours -------------------------------------------------------
\definecolor{kred}{HTML}{D81E2E}
\definecolor{kgrey}{HTML}{595959}

\usepackage[colorlinks=true,urlcolor=kred,linkcolor=kred]{hyperref}

% ---- Worksheet macros (inlined) ------------------------------------
\newcommand{\wbsection}[1]{\par\vspace{11pt}\noindent
  {\large\bfseries\color{kred}#1}\par\vspace{2pt}
  {\color{kred!35}\hrule height 1pt}\par\vspace{6pt}}

\newcommand{\task}[1]{\par\vspace{4pt}\noindent
  \textcolor{kred}{\textbf{$\triangleright$}}~#1\par\vspace{2pt}}

% Answers are discussed verbally -- no writing rules, just a small gap.
\newcommand{\answerlines}[1]{\par\vspace{0.4em}}

\newcommand{\boxlabel}[2]{\par\vspace{3pt}\noindent
  \fbox{\parbox{\dimexpr\linewidth-2\fboxsep-2\fboxrule}{\textbf{#1:}\enspace #2}}\par\vspace{3pt}}
\newcommand{\evidence}[1]{\boxlabel{Evidence to capture}{#1}}
\newcommand{\note}[1]{\boxlabel{Note}{#1}}
\newcommand{\caution}[1]{\boxlabel{Caution}{#1}}

% ---- Aha box + no italics (added) ----------------------------------
\newcommand{\aha}[1]{\par\vspace{5pt}\noindent
  \setlength{\fboxrule}{1.2pt}%
  {\color{kred}\fbox{\normalcolor\parbox{\dimexpr\linewidth-2\fboxsep-2\fboxrule}{%
     {\bfseries\color{kred}AHA.}\enspace #1}}}%
  \setlength{\fboxrule}{0.4pt}\par\vspace{5pt}}
\renewcommand{\emph}[1]{\textup{#1}}

\newcommand{\cmd}[1]{\texttt{#1}}
\newcommand{\cbox}{$\square$\enspace}

% ---- Hint and solutions divider -- same definitions as in kristiania-workbook.sty ----
\newcommand{\hint}[1]{\par\vspace{1pt}{\footnotesize\color{kgrey}\textbf{Hint:}\enspace #1}\par\vspace{2pt}}
\newcommand{\solheader}{\par\vspace{9pt}\noindent{\color{kred}\rule{\linewidth}{1.2pt}}\par\vspace{4pt}}

\newcommand{\signoff}{\par\vspace{9pt}\noindent
  {\color{kred}\rule{\linewidth}{0.4pt}}\par\vspace{3pt}
  \noindent{\footnotesize\color{kgrey}\textbf{Progress check:} Discuss your answers in the group, then
  talk the teacher through them verbally at the next session only to track progress; nothing is handed
  in, signed or graded.}\par}

\setlist[enumerate]{itemsep=3pt,topsep=2pt,parsep=0pt}
\setlist[itemize]{itemsep=2pt,topsep=2pt,parsep=0pt}
\setlength{\parindent}{0pt}
\setlength{\parskip}{2pt}

\begin{document}
\begin{center}
  {\LARGE\bfseries Scan your app, then shield it with a WAF}\\[3pt]
  {\large Session 8 -- DevSecOps, WAF \& secure apps}\\[5pt]
  {\small Exercise 8 \textperiodcentered\ Session 8 lab \textperiodcentered\ Estimated time: 90--120 min}
\end{center}
\vspace{2pt}\hrule\vspace{1.1em}

Write your answers in the answer document \cmd{SKY2100\_Exercise8\_Answers.docx}. Your Nextcloud is encrypted (S6) and the host is hardened (S7). Now you treat it as
an \emph{application} to test and protect. You will run a dynamic (DAST) scan, check for outdated
components, then put a web application firewall in front and watch it block a test attack.

\noindent\textbf{Caution:} Scan and attack only your \emph{own} VM. Sending attack payloads at systems
you do not own is illegal. The test payloads below are harmless probes against your own service.

\wbsection{1\quad Dynamic scan (DAST)}
\task{Run an OWASP ZAP baseline scan against your HTTPS Nextcloud (Docker one-liner, or the ZAP GUI).}
Use the VM's address from \cmd{ip a}, not \cmd{localhost}: inside the container, \cmd{localhost} is the
container itself.
\begin{quote}\ttfamily
docker run --rm -t ghcr.io/zaproxy/zaproxy:stable \textbackslash\\
\quad zap-baseline.py -t https://<vm-ip>
\end{quote}
\task{Skim the report. Name two findings and their risk level (e.g.\ missing security header, outdated
component). You do not need to fix them all today.}

\wbsection{2\quad Check your components}
\task{Open Nextcloud's built-in \emph{Security \& setup warnings} (Administration settings $\rightarrow$
Overview) and note any warning or outdated component.}
\task{Why is an outdated dependency an application-security risk, even if \emph{your} code is fine?}

\wbsection{3\quad Put a WAF in front (ModSecurity + OWASP CRS)}
\task{Add ModSecurity with the OWASP Core Rule Set to your reverse proxy in front of Nextcloud.}
A quick route is the official ModSecurity-CRS container, or the Nginx/Apache ModSecurity module. Start
in \emph{detection} mode (\cmd{SecRuleEngine DetectionOnly}), confirm the site still works, then switch
to \emph{prevention} (\cmd{SecRuleEngine On}).
\task{Why start in detection mode before turning on prevention? (one line)}

\wbsection{4\quad Attack your own app --- and watch it block}
\task{Send a harmless test payload that the OWASP CRS should catch, e.g.:}
\begin{quote}\ttfamily
curl -k "https://<your-host>/?q=<script>alert(1)</script>"\\
curl -k "https://<your-host>/?id=1' OR '1'='1"
\end{quote}
\task{What HTTP status do you get with the WAF in prevention mode, and what was logged?}

\aha{Think about what your earlier controls did and did not cover. TLS in Session~6 encrypted the
connection, and the hardening in Session~7 locked down the host, yet the payload you just sent was a
perfectly valid HTTPS request to a port you meant to keep open. Encryption and a firewall cannot see
that the content is an attack, because to them it is normal traffic to an allowed service. That is why
the attack moves up to the application layer, and why the WAF works there: it reads the request itself
and blocks known bad patterns before they reach the app. Note also that the WAF matched a known
pattern, so it is a guard, not a cure. A novel logic flaw it has no rule for would pass, which is why
secure code and scanning still matter.}

\wbsection{5\quad Cloud transfer: Microsoft Learn}
\task{Read ``Tutorial: Create an Application Gateway with a Web Application Firewall'' and note how the
WAF uses the OWASP rule set.}
\par\noindent\mbox{\small\href{https://learn.microsoft.com/azure/web-application-firewall/ag/application-gateway-web-application-firewall-portal}{learn.microsoft.com/azure/web-application-f{}irewall/ag/application-gateway-web-application-f{}irewall-portal}}
\task{Azure WAF runs in \emph{detection} or \emph{prevention} mode. Which one actually protects the app,
and why would you ever run the other?}

\wbsection{6\quad Azure reflection}
\task{In a DevSecOps pipeline, \textbf{SAST} runs in CI and \textbf{DAST} runs in CD. In one or two
lines each, say what each test looks at and why both are needed.}
\task{Your WAF blocked an attack at runtime. Name one application-security control from \emph{earlier} in
the SSDLC that should also address the same attack (defence in depth).}

\wbsection{7\quad Knowledge check}
\begin{enumerate}
\item Why does ``the attack move up the stack'' to the application after S6 and S7?
\hint{Transit (S6) and the host (S7) are locked, so the application is what remains.}
\item Name three OWASP Top~10 categories.
\hint{Broken access control, injection, cross-site scripting — any three.}
\item What is the common root cause of injection and XSS?
\hint{Untrusted input treated as code or commands.}
\item Why is a vulnerable \textbf{dependency} your problem, even with perfect code?
\hint{Most of a modern app is third-party and open-source code.}
\item List the five stages of the CSA Secure SDLC.
\hint{Design, coding, build/integrate/test, deliver/deploy, operate.}
\item What is \textbf{threat modelling}, and when in the lifecycle is it done?
\hint{Finding and ranking threats — done early, during design.}
\item Explain ``shift left'' and why it saves money.
\hint{Move security toward the start; a design flaw is far cheaper than a production one.}
\item Define \textbf{SAST} and \textbf{DAST}, and say when each runs in CI/CD.
\hint{One reads source code before deployment; one probes the running app.}
\item What is \textbf{SCA}, and what is an \textbf{SBOM}?
\hint{It audits your external dependencies for known vulnerabilities; the SBOM is their ingredient list.}
\item What does a \textbf{WAF} do, and which OWASP resource backs its rules?
\hint{It inspects HTTP at layer 7; its rules come from the OWASP Core Rule Set.}
\item Detection vs prevention mode: what is the difference?
\hint{One logs matching requests but lets them through; one blocks them.}
\item Why is a WAF ``a guard, not a cure''? Give one limitation.
\hint{It blocks known patterns — a novel logic flaw walks straight past.}
\end{enumerate}

\signoff

\clearpage
\solheader
{\LARGE\bfseries\color{kred} Solutions}\\[2pt]
{\small Work through the lab first, then check here. Model answers are indicative — equivalent reasoning is fine.}\par\vspace{8pt}
\noindent\textbf{Note:} Model answers are indicative -- accept equivalent reasoning. Multiple-choice
answers are in \textbf{bold}.

\wbsection{Knowledge check}
\begin{enumerate}
\item Each control raised the bar: TLS protects transit (S6) and hardening locks the host (S7), so the
      remaining target is the \emph{application logic}. A valid HTTPS request to an open service can
      still carry an attack in its payload, which the firewall lets through.
\item Any three of: broken access control; injection (SQL/command); cross-site scripting (XSS); security
      misconfiguration; vulnerable/outdated components; identification \& authentication failures;
      cryptographic failures; insecure design.
\item \textbf{Untrusted input treated as code/commands} -- a SQL string, a shell command, or browser
      script. Fix by validating input and parameterising/encoding output.
\item Modern apps are mostly third-party and open-source code; a flaw in a dependency is a flaw in your
      app (e.g.\ Log4Shell). You ship the dependency, so you own its risk -- hence SCA and an SBOM.
\item \textbf{Secure Design \& Architecture}; \textbf{Secure Coding}; \textbf{Continuous Build,
      Integration \& Testing}; \textbf{Continuous Delivery \& Deployment}; \textbf{Runtime Defense \&
      Monitoring} (CSA p.~231).
\item A structured process to identify, assess and prioritise threats to a system's objectives. Done
      \emph{early}, during design, when changing the answer is cheap (CSA p.~231--232).
\item ``Shift left'' moves security toward the start of the lifecycle. A flaw caught in design/code is
      far cheaper to fix than one found in production (CSA p.~233).
\item \textbf{SAST} reads source code at rest for flaws -- runs pre-deployment in \textbf{CI}.
      \textbf{DAST} probes the running app as a black box -- runs post-deployment in \textbf{CD}. Both
      are needed: code vs runtime (CSA p.~234--235, 242).
\item \textbf{SCA (Software Composition Analysis)} audits external/open-source dependencies for known
      vulnerabilities and licensing risk. An \textbf{SBOM (Software Bill of Materials)} lists all
      components shipped, giving transparency (CSA p.~234).
\item A \textbf{WAF} inspects HTTP requests at layer~7 and blocks known attack patterns (SQLi, XSS)
      before they reach the app. Its rules are based on the \textbf{OWASP Core Rule Set (CRS)}.
\item \textbf{Detection} logs matching requests but lets them through (for tuning); \textbf{prevention}
      actively blocks them. Only prevention protects the app.
\item It blocks \emph{known} patterns and can be bypassed by novel logic flaws; it needs tuning (too
      strict $\rightarrow$ false positives, too loose $\rightarrow$ misses). It complements secure code,
      it does not replace it.
\end{enumerate}

\wbsection{Lab checkpoints}
\begin{itemize}
\item \textbf{DAST (ZAP):} the baseline scan typically flags missing security headers (CSP, HSTS),
      cookie flags, and sometimes outdated components. The point is to \emph{read} a dynamic report --
      DAST is the post-deployment, black-box test (CSA p.~234).
\item \textbf{Component check:} Nextcloud's security/setup warnings stand in for SCA -- outdated PHP,
      missing modules, or an out-of-date instance. The lesson: dependencies and configuration are part
      of the attack surface.
\item \textbf{WAF in front:} ModSecurity + OWASP CRS should sit between the client and Nextcloud.
      Starting in \emph{DetectionOnly} avoids blocking legitimate traffic while tuning; then
      \emph{SecRuleEngine On} enforces. This is the SSDLC runtime-defence stage.
\item \textbf{Blocked attack:} the XSS/SQLi probes should return \textbf{403} in prevention mode, with a
      matching CRS rule ID in the audit log (e.g.\ a 941xxx XSS or 942xxx SQLi rule). Detection mode
      would log but still return 200 -- a good contrast to show.
\end{itemize}

\wbsection{Cloud transfer \& reflection}
\begin{itemize}
\item \textbf{SAST vs DAST:} SAST reads source code for flaws/secrets in CI (pre-deployment); DAST
      probes the running app from the outside in CD (post-deployment). Both are needed because each has
      the other's blind spot -- code internals vs real runtime behaviour (CSA p.~242).
\item \textbf{Defence in depth:} the WAF blocks the injection/XSS at runtime, but the \emph{same} attack
      should also be addressed earlier -- secure coding (input validation, parameterised queries, output
      encoding) and SAST in CI. Fix it in the code \emph{and} block it at the edge.
\item \textbf{Detection vs prevention:} prevention blocks and protects; detection only logs. You run
      detection first to tune out false positives before enforcing, so real users are not blocked.
\end{itemize}

\wbsection{References}
CSA Security Guidance v5, \textbf{Domain~10: Application Security} (p.~230--245, verified): SSDLC five
stages (p.~231), threat modelling (p.~231--232), SAST/SCA/static scanning (p.~233--234), DAST/fuzzing/
IAST/pen test/bug bounty (p.~234--235), DevSecOps CI$\rightarrow$SAST / CD$\rightarrow$DAST and the six
pillars (p.~242--244). Comer, \emph{The Cloud Computing Book}, Ch.~15 -- DevOps: the DevOps approach
(p.~210), CI (p.~211), CD (p.~212), sandbox/canary/blue-green (p.~212--213), difficult aspects
(p.~213--214), verified. MS Learn: \emph{Application Gateway with a Web Application Firewall} (verified,
June 2026) -- OWASP CRS, detection vs prevention mode.

\end{document}
